\section*{Capítulo \ref{ch:g12:synthesis-strategy} --- Estrategia de síntesis y química verde}

\begin{solution}{exo:g12:synthesis-strategy:1}
Hidratación: $46.0/(28.0 + 18.0) = \qty{100}{\percent}$. Fermentación:
$2 \times 46.0/180.0 = \qty{51.1}{\percent}$.
\end{solution}

\begin{solution}{exo:g12:synthesis-strategy:2}
Los dos grupos, \ce{Z} sobre la \omterm{def:g11:functional-groups:amine}{amina} de la glicina y el \omterm{def:g11:functional-groups:carboxylic-acid}{éster} metílico
sobre el ácido de la alanina, se colocan en la etapa 1 y se retiran en
la etapa 3.
\end{solution}

\begin{solution}{exo:g12:synthesis-strategy:3}
El principio 5: usar \omterm{def:g6:solutions-and-solubility:solute}{disolventes} y condiciones de reacción más seguros.
\end{solution}

\begin{solution}{exo:g12:synthesis-strategy:4}
Un exceso de etanol hace que $Q < K$, así que la reacción avanza más en
sentido directo. Un \omterm{def:g12:catalysis:catalyst}{catalizador} acelera por igual los dos sentidos: se
alcanza el mismo \omterm{def:g10:reaction-progress-table:system}{estado final}, solo que antes.
\end{solution}

\begin{solution}{exo:g12:synthesis-strategy:5}
Sí: solo se reduce la \omterm{def:g11:functional-groups:carbonyl}{cetona}, a un grupo \omterm{def:g11:functional-groups:alcohol}{alcohol} secundario; el \omterm{def:g11:functional-groups:carboxylic-acid}{éster} no
cambia.
\end{solution}

\begin{solution}{exo:g12:synthesis-strategy:6}
Más \omterm{def:g11:functional-groups:carboxylic-acid}{éster}: un exceso de uno de los \omterm{def:g7:chemical-reactions:reactant}{reactivos}, o retirar el agua (trampa
de agua) a medida que se forma. Más deprisa: calentar a reflujo y añadir
un \omterm{def:g12:catalysis:catalyst}{catalizador} ácido.
\end{solution}

\begin{solution}{exo:g12:synthesis-strategy:7}
Con \ce{NaBH4}:
\[ \ce{CH3-CH(OH)-CH2-CH2-CO-O-CH3} . \]
Con \ce{LiAlH4}:
\[ \ce{CH3-CH(OH)-CH2-CH2-CH2OH} \]
(pentano-1,4-diol), más metanol procedente de la otra mitad del \omterm{def:g11:functional-groups:carboxylic-acid}{éster}.
\end{solution}

\begin{solution}{exo:g12:synthesis-strategy:8}
La adición (\qty{100}{\percent}), la esterificación
(\qty{83}{\percent}) y la ruta de tres etapas del ibuprofeno
(\qty{77}{\percent}). $77.4/40.0 = 1.9$: casi el doble de buena.
\end{solution}

\begin{solution}{exo:g12:synthesis-strategy:9}
$180.0/(138.0 + 102.0) = \qty{75.0}{\percent}$. Ácido etanoico:
$60.0/180.0 = \qty{0.333}{kg}$ por kilogramo de aspirina.
\end{solution}

\begin{solution}{exo:g12:synthesis-strategy:10}
El calentamiento acelera la reacción; el refrigerante devuelve los
vapores al matraz, de modo que no se pierde \omterm{def:g7:chemical-reactions:reactant}{reactivo} ni \omterm{def:g6:solutions-and-solubility:solute}{disolvente}. El
calentamiento mejora la velocidad, no el \omterm{def:g10:synthesis-yield:yield}{rendimiento} final.
\end{solution}

\begin{solution}{exo:g12:synthesis-strategy:11}
Seis etapas frente a tres. Principios: prevenir los residuos (1),
maximizar la \omterm{def:g12:synthesis-strategy:atom-economy}{economía atómica} (2), usar \omterm{def:g12:catalysis:catalyst}{catalizadores} (9); además, menos
etapas, menos energía (6).
\end{solution}

\begin{solution}{exo:g12:synthesis-strategy:12}
$0.80 \times 0.70 \times 0.90 = 0.504$: \qty{50.4}{\percent}. Material de
partida: $0.10/0.504 = \qty{0.20}{mol}$.
\end{solution}

\begin{solution}{exo:g12:synthesis-strategy:13}
Proteger la \omterm{def:g11:functional-groups:amine}{amina} de la alanina (\ce{Z}) y el ácido de la glicina (\omterm{def:g11:functional-groups:carboxylic-acid}{éster}
metílico); formar el enlace \omterm{def:g11:functional-groups:amine}{amida} entre el ácido libre de la alanina y la
\omterm{def:g11:functional-groups:amine}{amina} libre de la glicina; eliminar los dos \omterm{def:g12:synthesis-strategy:protecting-group}{grupos protectores}. Tres
fases, es decir, cinco reacciones si se cuenta cada protección y cada
desprotección.
\end{solution}

\begin{solution}{exo:g12:synthesis-strategy:14}
El \omterm{def:g12:catalysis:catalyst}{catalizador} aplica el principio 9. Pero \qty{250}{\celsius} va contra
el principio 6, el benceno contra los principios 3 y 5, y una \omterm{def:g12:synthesis-strategy:atom-economy}{economía
atómica} del \qty{35}{\percent} contra el principio 2: $65/35 = \qty{1.9}{kg}$
de residuos por kilogramo de producto. Un solo rasgo verde no hace verde
un proceso.
\end{solution}

\begin{solution}{exo:g12:synthesis-strategy:15}
% ledger: crit:CO2-T, crit:CO2-p
\qty{40}{\celsius} y \qty{100}{bar} están por encima del punto crítico
del dióxido de carbono, \qty{31}{\celsius} y \qty{73.8}{bar}. La presión
es cien veces la presión atmosférica: el recipiente es de acero grueso.
El dióxido de carbono no es tóxico ni inflamable, vuelve a ser un gas sin
dejar residuo y puede reutilizarse; no se escapa ningún \omterm{def:g6:solutions-and-solubility:solute}{disolvente}
clorado.
\end{solution}

\begin{solution}{pb:g12:synthesis-strategy:1}
% ledger: ibu:bhc-ae-recovered
\textbf{1.} C: $10 + 4 + 1 = 15 = 13 + 2$; H: $14 + 6 + 2 = 22 = 18 + 4$;
O: $3 + 1 = 4 = 2 + 2$.

\textbf{2.} El ácido etanoico, \ce{C2H4O2}.

\textbf{3.} El nitrógeno, el cloro y el sodio: el ibuprofeno solo
contiene carbono, hidrógeno y oxígeno.

\textbf{4.} La nueva: principio 9, \omterm{def:g12:catalysis:catalyst}{catalizadores} en lugar de \omterm{def:g7:chemical-reactions:reactant}{reactivos}
estequiométricos.

\textbf{5.} $13 \times 12.0 + 18 \times 1.0 + 2 \times 16.0 =
\qty{206.0}{g/mol}$.

\textbf{6.} $134.0 + 102.0 + 122.5 + 68.0 + 19.0 + 33.0 + 36.0 =
\qty{514.5}{g/mol}$.

\textbf{7.} $206.0/514.5 = \qty{40.0}{\percent}$.

\textbf{8.} $134.0 + 102.0 + 2.0 + 28.0 = \qty{266.0}{g/mol}$.

\textbf{9.} $206.0/266.0 = \qty{77.4}{\percent}$.

\textbf{10.} $(206.0 + 60.0)/266.0 = \qty{100}{\percent}$ sobre el
papel; más del \qty{99}{\percent} en la práctica.

\textbf{11.} $514.5/206.0 = \qty{2.50}{t}$.

\textbf{12.} $2.50 - 1 = \qty{1.50}{t}$.

\textbf{13.} $266.0/206.0 = \qty{1.29}{t}$.

\textbf{14.} $60.0/206.0 = \qty{0.29}{t}$.

\textbf{15.} $1.50 - 0.29 = \qty{1.21}{t}$.

\textbf{16.} $0.90^6 = 0.53$: \qty{53}{\percent}.

\textbf{17.} $0.90^3 = 0.73$: \qty{73}{\percent}.

\textbf{18.} Cada etapa usa \omterm{def:g6:solutions-and-solubility:solute}{disolventes}, separaciones y purificaciones, y
pierde algo de producto: menos etapas, menos pérdidas de este tipo.

\textbf{19.} Casi ninguna: su único subproducto se aprovecha.

\textbf{20.} La ruta de tres etapas evita unas \qty{1.5}{t} de residuos
por tonelada de ibuprofeno (\qty{1.2}{t} incluso si su ácido etanoico se
tirara).
\end{solution}
